% ---------------------------------------------------------------------------
% Results fragment for Section 10 (allied network design), Results subsection
% (subsec:network_results), and for Appendix C (calibration and robustness).
% Final design: 2026-10-12 full rerun, scripts/30 at lambda = 0.78, 37 network nodes (the EU-27 as
% one node), the 22 Upstream and Midstream stages, GDP-log portfolio caps on, node policy term off.
% Floats (label = file name). The placeholders of 10_network.tex map as follows:
%   fig:specialization -> fig:network_specialization   fig:concentration -> fig:network_concentration
%   fig:portfolio      -> fig:network_portfolio        fig:network_map   (unchanged)
%   new: fig:network_frontier, fig:network_robustness, tab:network_producers, tab:network_eu_split,
%        tab:network_frontier, tab:network_robustness (tables in paper/tables/, \input below)
% Numbers: paper/results/network.json (scripts/paper/fig_network.R),
% paper/results/network_frontier.json (scripts/paper/fig_network_frontier.R) and
% paper/results/network_robustness.json (scripts/paper/fig_network_robustness.R).
% Citations still to be added to references.bib are marked "% CITE:".
% ---------------------------------------------------------------------------

\paragraph{Coverage and solution.} The production design covers the 57 coalition members (31 core
members and 26 partners) and the six EU member states that belong to the coalition only through
the EU node (Cyprus, Estonia, Latvia, Lithuania, Malta and Slovenia): 63 economies, which form
37 network nodes once the EU-27 is bundled into one. It solves the 22 Upstream and Midstream
technology--stage programs; Downstream deployment is not part of the allied division of labour.
Every node enters every stage with a producer score, which since the policy term was switched off
rests on $\EO$ and $\ES$ alone. In two stages, Geothermal Upstream and Midstream, no node has an
$\ES$ value, so all 37 producer scores rest on $\EO$; in Green Hydrogen Midstream five small
partners (Costa Rica, the Dominican Republic, Ghana, Guatemala and Namibia) lack $\ES$. Without
an $\ES$ value demand rests on size alone, so Namibia, the smallest node, has zero demand in these
three stages (its S-curved size is zero); it remains a potential producer in them. The inputs arrive with one record per key: for each of the five
collapsed inputs (2{,}079 $\EO$ and 1{,}881 $\ES$ records, 726 development-potential records and
128{,}898 friendshore and 116{,}622 opportunity dyads among the 63 economies) the number of rows
read equals the number of keys kept, so the collapse rules of Section~\ref{subsec:network_sets}
act as guards only. All 22 programs are solved by the MILP and none falls back to the greedy
heuristic. The tangent approximation of the Herfindahl term, on a share grid of 0.01, has a
largest error of $2.5\times10^{-5}$ per producer, its theoretical bound. The portfolio loop
converges after seven passes with no violation remaining. On a shared desktop workstation
(R~4.4.1, GLPK through \texttt{ompr} and \texttt{ROI}) the production design takes 462 seconds
of wall-clock time for 62 stage solves, 22 in the first pass and 40 re-solves in the portfolio
loop; a first-pass stage solve takes 5.9 to 8.6 seconds (median 7.2, slowest Wind Midstream) and
no solve exceeds 10.6 seconds, far inside the 300-second stage limit.

\paragraph{Who produces what.} Figure~\ref{fig:network_specialization} and
Table~\ref{tab:network_producers} show the designed division of labour. The 22 stages have 201
producer roles, between seven producers (Nuclear Upstream) and ten (eight stages), with a median
of nine in both Upstream and Midstream. The Herfindahl index of designed production shares,
$\HHI^{\mathrm{D}}_{ts}$, ranges from 0.126 (Oil Midstream) to 0.209 (Nuclear Upstream) with a
median of 0.152, or 6.58 equally sized producers (4.78 to 7.96). The largest single share is
0.325 (Indonesia in Green Hydrogen Upstream and Canada in Nuclear Upstream), the largest share of
a stage is at most 0.30 in 18 of the 22 stages, and no share reaches the cap $\bar s=0.40$;
61 of the 201 shares (30\%) sit at the floor $\underline s=0.05$. This is the work of the
diversification term. In the linear design ($\lambda=0$) two producers sit at the 0.40 cap in
most stages and the median stage has three producers and an index of 0.345
(Figure~\ref{fig:network_frontier}); the concave penalty moves the solution off these corners, so
shares are interior and set jointly by the node scores, the edges and the penalty.

Twenty-nine of the 37 nodes produce somewhere; eight never do (Turkey, New Zealand, Ukraine, the
Dominican Republic, Guatemala, Ghana, Costa Rica and Gabon). Canada has the most roles (16, nine
of them Upstream), followed by the EU node (15), the United States and Indonesia (13 each), India
and Japan (12 each), and Thailand and Mexico (11 each). Eight nodes produce only Upstream, all of
them resource economies (Australia, Chile, Saudi Arabia, Peru, Argentina, Oman, Namibia and the
Democratic Republic of the Congo); Japan takes ten of its twelve roles and Korea nine of its ten
in Midstream. Across nodes the Spearman rank correlation between the number of roles and the
portfolio cap is 0.75 (0.74 with GDP), so the breadth of a node's portfolio is largely a matter
of size, through the caps, while which stages it takes follows its scores. The 22 nodes outside
the high-income group carry 47\% of designed Upstream production and 43\% of Midstream
production, and the partner tier 58\% and 45\%.

Three patterns in the producer map arise from the construction rather than from observed
capacity, and the reader should weigh them accordingly. First, Indonesia leads six Upstream
stages (Batteries, Coal, Electric Vehicles, Green Hydrogen, Solar and Wind, with shares of 0.225
to 0.325) and Oil Midstream. Indonesia is China-leaning by its UN votes and enters the coalition
through IPEF (Section~\ref{sec:coalition}); the producer score carries no alignment term, and
alignment reaches the design only through the friendshore edges, so a partner with strong $\EO$
and $\ES$ is selected on capability, whatever its alignment. Second, Thailand leads both
Geothermal stages, in which the producer score rests on $\EO$ alone and, as shown below, the
edges lack their opportunity component: these are the stages in which the design rests on the
least information. Third, the design applies no endowment screen of the kind used for the
allied picks of the Strategic Index (Section~\ref{sec:strategic}), so a node can be a designed
Upstream producer without recorded reserves or production of the stage's resources. Applying the
screen's rule after the fact (technology-level reserves or production above zero; the EU node
counts as endowed when any member is), 21 of the 90 designed roles in the ten screened Upstream
stages (all but Geothermal) are held by nodes without a recorded resource base. They carry a
median of 11\% of a screened stage's production: none in Electric Vehicles, Gas and Oil Upstream,
5\% in Batteries Upstream (the United Kingdom), 10\% in Nuclear Upstream (Israel and Thailand at
the floor), 12.5\% in Coal Upstream (Peru and the Philippines), 23.5\% in Wind Upstream (Japan,
Israel, Korea and Singapore), 27.5\% in Electric Grid Upstream (five nodes at or near the floor),
28\% in Green Hydrogen Upstream (Thailand 0.145, Israel and the United Kingdom) and 36\% in Solar
Upstream (the United Kingdom, Norway, South Africa and Singapore). These nodes score on the stage
through its trade basket, readiness and the other categories of $\EO$ and $\ES$, not through a
domestic resource. A reader who equates Upstream production with extraction should read their
shares as supply that the node would have to source through trade or overseas equity; the screen
is a rule for the Strategic Index picks, and we do not impose it on the design.

% Generated by scripts/paper/fig_network.R from data/processed/outputs/allied_network_specialization.csv
\begin{table}[htbp]
\centering
\begin{threeparttable}
\caption{Designed producers and production shares by technology and stage}
\label{tab:network_producers}
\footnotesize
\begin{tabularx}{\textwidth}{@{}>{\raggedright\arraybackslash}p{2.1cm} >{\raggedright\arraybackslash}X >{\raggedright\arraybackslash}X@{}}
\toprule
Technology & Upstream & Midstream \\
\midrule
Batteries & IDN~.23, AUS~.22, CAN~.12, CHL~.12, ARG~.08, PHL~.06, EU~.06, GBR~.05, ZAF~.05 & USA~.29, IND~.14, MYS~.14, CAN~.11, EU~.10, THA~.06, KOR~.06, GBR~.05, JPN~.05 \\
Coal & IDN~.26, IND~.15, VNM~.12, COL~.09, ZAF~.07, PER~.07, CAN~.06, THA~.06, EU~.05, PHL~.05 & JPN~.20, EU~.20, COL~.14, KOR~.12, VNM~.08, CAN~.07, ISR~.06, MEX~.05, MYS~.05 \\
Electric Grid & CHL~.31, MEX~.14, COD~.12, PER~.12, PHL~.07, AUS~.05, GBR~.05, JPN~.05, NOR~.05, ZAF~.05 & THA~.22, MEX~.15, JPN~.13, MYS~.11, EU~.10, USA~.08, CAN~.05, IDN~.05, IND~.05, MAR~.05 \\
Electric Vehicles & IDN~.24, AUS~.21, USA~.16, CHL~.08, PHL~.07, ARG~.06, EU~.05, THA~.05, ZAF~.05 & USA~.32, MEX~.20, KOR~.19, JPN~.08, VNM~.06, CAN~.05, GBR~.05, ZAF~.05 \\
Gas & USA~.25, CAN~.14, NOR~.14, QAT~.12, MYS~.10, IDN~.05, AUS~.05, IND~.05, OMN~.05, SAU~.05 & IND~.24, USA~.18, EU~.14, KOR~.12, MEX~.06, GBR~.05, IDN~.05, ISR~.05, JPN~.05, NOR~.05 \\
Geothermal & THA~.26, MYS~.15, NOR~.13, ISR~.10, EU~.08, IDN~.07, MEX~.05, CAN~.05, QAT~.05, USA~.05 & THA~.26, MYS~.15, ISR~.14, NOR~.10, QAT~.07, EU~.07, CAN~.06, IDN~.05, MEX~.05, USA~.05 \\
Green Hydrogen & IDN~.33, THA~.14, CAN~.14, EU~.12, MEX~.09, ISR~.08, GBR~.05, PHL~.05 & USA~.26, EU~.21, VNM~.14, NOR~.08, PHL~.07, SGP~.07, KOR~.06, ISR~.05, JPN~.05 \\
Nuclear & CAN~.33, NAM~.22, IND~.16, ZAF~.14, ISR~.05, SAU~.05, THA~.05 & IND~.24, EU~.20, JPN~.14, MYS~.12, ISR~.07, KOR~.06, NOR~.06, GBR~.05, MEX~.05 \\
Oil & USA~.23, SAU~.17, CAN~.17, NOR~.12, OMN~.08, QAT~.08, IDN~.05, IND~.05, THA~.05 & IDN~.20, IND~.14, KOR~.14, THA~.13, VNM~.09, JPN~.08, EU~.06, CAN~.05, MYS~.05, NOR~.05 \\
Solar & IDN~.24, CAN~.15, EU~.12, GBR~.12, NOR~.10, ZAF~.10, IND~.07, SGP~.06, MEX~.05 & IND~.28, USA~.20, VNM~.14, THA~.12, PHL~.06, CAN~.05, JPN~.05, KOR~.05, SGP~.05 \\
Wind & IDN~.22, USA~.22, AUS~.20, JPN~.08, ZAF~.06, CAN~.06, ISR~.05, KOR~.05, SGP~.05 & IND~.28, USA~.25, EU~.15, JPN~.09, MEX~.06, BRA~.06, KOR~.05, VNM~.05 \\
\bottomrule
\end{tabularx}
\begin{tablenotes}[flushleft]\footnotesize
\item \emph{Notes:} Producers of each stage in the production design (37 network nodes, diversification weight $\lambda=0.78$, portfolio caps on, every stage solved by the MILP), ordered by the designed production share $x_{p,ts}$, the share of the coalition's unit demand in the stage that the node supplies (ISO3 codes; EU = the EU-27 node; shares sum to one within a cell and are printed to two decimals). Source: \texttt{allied\_network\_specialization.csv}.
\end{tablenotes}
\end{threeparttable}
\end{table}


\begin{figure}[htbp]
  \centering
  \includegraphics[width=\textwidth]{network_specialization}
  \caption{Designed division of labour across the allied network. Rows are the 37 network nodes
  (the EU-27 as one node), ordered by the number of stages in which they produce (in
  parentheses) and then by their summed production share; columns are the eleven headline
  technologies within the Upstream and Midstream stages. Shading is the designed production
  share $x_{p,ts}$, the share of the network's unit demand in the stage that the node supplies,
  on a sequential scale from the floor 0.05 to the largest share, 0.325; light grey cells are
  stages in which the node does not produce. A cross marks a stage from which the portfolio-cap
  loop barred the node in its final pass. Production design: 37 nodes, diversification weight
  $\lambda=0.78$, GDP-log portfolio caps, node policy term off, all 22 stages solved by the MILP.
  Source: \texttt{allied\_network\_specialization.csv} and
  \texttt{allied\_network\_diagnostics.csv}.}
  \label{fig:network_specialization}
\end{figure}

\paragraph{The EU node.} The EU node produces in 15 of the 22 stages (six Upstream, nine
Midstream), with shares from 0.05 to 0.215 and a median of 0.105. It is the second-largest
producer in Green Hydrogen Midstream (0.215) and Nuclear Midstream (0.205), and its 0.205 in Coal
Midstream equals Japan's leading share, but it is the sole largest producer in no stage. It is a top-5 producer in ten stages, exactly its cap of ten, and the
portfolio loop bars it from five (Electric Grid Upstream, Electric Vehicles Midstream, Nuclear
Upstream, Solar Midstream and Wind Upstream), so Electric Vehicles and Solar are the two
Midstream stages without EU production. By its observed extra-EU exports the EU is the largest
exporting node in nine of the 18 stages with a trade basket, all of them Midstream. The design
sees the EU differently from the trade data in two ways. Its producer score is the stage-weighted
mean of 27 members' $\EO$ and $\ES$, so strong and weak members are averaged, and intra-EU trade,
which the extra-EU benchmark leaves out, is internal to the node. And its demand is
rank-compressed: demand mixes S-curved ranks of need and size, so the EU, with 24.6\% of the
nodes' GDP, has a median demand weight of 0.035 across stages (0.029 to 0.046), the third largest
of the nodes, against a largest weight of any node in any stage of 0.054; the United States, with
35.6\% of the nodes' GDP, has 0.033. In the design the EU therefore consumes about as much as
the United States and not much more than a mid-sized node.
Table~\ref{tab:network_eu_split} splits the EU's designed shares back to its members in
proportion to the stage weights that aggregate their inputs; summed over stages, Germany (0.41),
France (0.24), Poland (0.19) and Italy (0.16) carry the largest derived shares. These member
figures are derived, not designed: the program never sees the members separately.

% Generated by scripts/paper/fig_network.R from data/processed/outputs/allied_network_member_shares.csv
\begin{table}[htbp]
\centering
\begin{threeparttable}
\caption{The EU node's designed production shares and their split to member states}
\label{tab:network_eu_split}
\footnotesize
\begin{tabular}{@{}l l r l l r@{}}
\toprule
Technology & Stage & EU share & Weights & Three largest members (derived share) & Top-3 weight \\
\midrule
Green Hydrogen & Upstream & .125 & GDP & DEU~.030, FRA~.020, ITA~.015 & .52 \\
Solar & Upstream & .115 & exports & FRA~.053, DEU~.020, NLD~.012 & .74 \\
Geothermal & Upstream & .085 & GDP & DEU~.020, FRA~.013, ITA~.010 & .52 \\
Batteries & Upstream & .055 & exports & DEU~.016, FRA~.006, FIN~.005 & .49 \\
Coal & Upstream & .050 & production & POL~.025, DEU~.014, CZE~.006 & .91 \\
Electric Vehicles & Upstream & .050 & exports & FIN~.014, DEU~.007, NLD~.007 & .54 \\
Green Hydrogen & Midstream & .215 & GDP & DEU~.051, FRA~.034, ITA~.026 & .52 \\
Nuclear & Midstream & .205 & exports & FRA~.058, DEU~.054, NLD~.028 & .68 \\
Coal & Midstream & .205 & exports & POL~.082, DEU~.032, ITA~.022 & .66 \\
Wind & Midstream & .155 & exports & DEU~.045, ITA~.016, POL~.012 & .47 \\
Gas & Midstream & .135 & exports & DEU~.038, ITA~.025, FRA~.010 & .54 \\
Batteries & Midstream & .105 & exports & POL~.023, HUN~.022, DEU~.020 & .62 \\
Electric Grid & Midstream & .095 & exports & DEU~.031, ITA~.008, FRA~.006 & .47 \\
Geothermal & Midstream & .070 & GDP & DEU~.017, FRA~.011, ITA~.008 & .52 \\
Oil & Midstream & .055 & exports & DEU~.016, ITA~.008, FRA~.004 & .51 \\
\bottomrule
\end{tabular}
\begin{tablenotes}[flushleft]\footnotesize
\item \emph{Notes:} Stages in which the EU-27 node produces in the production design (15 of 22). EU share: the node's designed production share. The design solves the EU as one node; its share is split back to the 27 member states in proportion to the stage weights that aggregate their inputs (Weights: the members' production or export shares in the stage, else their GDP shares), so the member figures are derived, not designed. Top-3 weight: the combined weight of the three largest members within the node. ISO3 codes. Source: \texttt{allied\_network\_member\_shares.csv}.
\end{tablenotes}
\end{threeparttable}
\end{table}


\paragraph{Designed flows.} Figure~\ref{fig:network_map} shows the flows of three Midstream
stages. In Batteries Midstream the United States carries 0.295 of production, India and Malaysia
0.135 each, Canada 0.114 and the EU 0.105, with four more nodes near the floor; in Solar
Midstream India (0.275), the United States (0.205), Viet Nam (0.135) and Thailand (0.125) lead;
in Wind Midstream India (0.275), the United States (0.250) and the EU (0.155). Across all 22
stages the redundancy constraints bind exactly. Of the 811 consumer--stage pairs, 793 draw on
exactly $m_S=2$ suppliers and 18 on three, and in every stage the weakest active link carries
exactly $\varepsilon=10\%$ of its consumer's demand. Self-supply is negligible: it is positive in
two stages only (0.8\% of demand in Electric Grid Upstream and 0.2\% in Electric Vehicles
Midstream), because a node's own edge is worth the median edge \eqref{eq:net_self_edge} while the
edges to its best partners lie above it. Routing is not regional: in the three mapped stages 15
to 28\% of the flow stays within a World Bank region.

Routing follows the strongest partnerships. Over flows between distinct nodes the flow-weighted
mean edge weight is 0.874 in the median stage (0.385 to 0.910). The low end is Geothermal: every
ordered node pair has the friendshore component in all 22 stages, but the opportunity component
is complete in 19 stages only, covers 86.5\% of the pairs in Green Hydrogen Midstream (the five
nodes without $\ES$ have no opportunity dyads as consumers) and is absent in both Geothermal
stages. A missing component enters the edge weight as zero, so the Geothermal edges enter at no
more than half weight and their flow-weighted weight is 0.385 to 0.387, against 0.73 to 0.75 in
the two Green Hydrogen stages and 0.85 to 0.91 elsewhere. Across the 22 stages,
90.3\% of the 29{,}304 stage-specific ordered pairs of distinct nodes carry both
components and 9.7\% the friendshore component only (the two Geothermal stages and the five Green
Hydrogen Midstream consumers without $\ES$); no pair lacks both. The design routes its flows to
the stronger pairs: the unweighted mean edge weight over these pairs is 0.58 in the median stage,
against the flow-weighted 0.874, and the flow-weighted mean is the higher of the two in all 22
stages. The EU node's edges aggregate its members' edges, weighted by the members' 2024
basket trade with the partner where it exists (76\% of the EU edges) and by the stage weights
otherwise; edges between two EU members are internal to the node and dropped.

\begin{figure}[htbp]
  \centering
  \includegraphics[width=\textwidth]{network_map}
  \caption{Designed flows in three Midstream stages. In each panel, producers sit on the left,
  sized by their production share $x_p$ (printed next to the code; EU = the EU-27 node), and
  every one of the 37 nodes sits on the right as a consumer, in the same order in every panel:
  by World Bank region (shaded bands, labelled at the right) and by GDP within region. A line is a
  flow $f_{pc}$ from producer $p$ to consumer $c$, with width increasing in the flow measured as
  a share of the stage's unit demand. Consumers that also produce in the stage are in bold.
  Producers are placed at the flow-weighted mean position of their consumers (the barycentre step
  of a two-layer layout) so that few lines cross. Panel titles give the number of producers and
  $\HHI^{\mathrm{D}}_{ts}$. Self-supply is zero in all three stages. Production design
  ($\lambda=0.78$, portfolio caps, MILP). Source: \texttt{allied\_network\_flows.csv} and
  \texttt{allied\_network\_specialization.csv}; layout computed with \texttt{igraph}.}
  % CITE: csardi2006igraph
  \label{fig:network_map}
\end{figure}

\paragraph{Concentration.} Figure~\ref{fig:network_concentration} compares the designed
concentration with two observed benchmarks in the 18 stages that have an HS6 trade basket (all
but the Geothermal and Green Hydrogen stages): the export Herfindahl index of the 37 nodes
themselves, $\HHI^{\mathrm{N}}_{ts}$, from the reconciled 2024 Comtrade matrix with the EU
counting only its extra-EU exports, and that of world exports, $\HHI^{\mathrm{W}}_{ts}$, over all
reporters \citep{uncomtrade}. The calibration of $\lambda$ targets the median of the first
(0.1556, or 6.43 equivalent exporters), and the design meets it in the median by construction
(0.1520, 6.58 equivalent producers). Stage by stage, however, the designed and the observed
concentration are unrelated: the design is less concentrated than the nodes' exports in 10 of
the 18 stages and more concentrated in eight, and the Spearman correlation of
$\HHI^{\mathrm{D}}_{ts}$ with $\HHI^{\mathrm{N}}_{ts}$ across the 18 stages is $-0.07$. The
design compresses the spread of concentration: $\HHI^{\mathrm{N}}_{ts}$ runs from 0.095 to 0.355,
$\HHI^{\mathrm{D}}_{ts}$ from 0.126 to 0.209. Where the nodes' exports are most concentrated, in
Nuclear Upstream (Canada 54\% of the nodes' exports), Coal Upstream (Australia 41\%) and Nuclear
Midstream (the EU 48\%), the design spreads production; where exports are already dispersed, as
in Batteries Midstream ($\HHI^{\mathrm{N}}=0.095$), it concentrates it on the highest-scoring
producers. The largest designed share is below the largest node's export share in 13 of the 18
stages, and the designed and observed leaders coincide in four. Measured against world exports
the design is more concentrated in 13 of the 18 stages (median $\HHI^{\mathrm{W}}=0.101$, 9.9
equivalent exporters), but so are the nodes' own exports in 16: a closed network of 37 nodes
cannot match the dispersion of a world market with up to 139 reporters.

What the network changes is the identity of the suppliers. China is the largest world exporter
in ten of the 18 basket stages: eight of the nine Midstream stages with a basket (Nuclear
Midstream, led by France, is the exception) and the Solar and Wind Upstream baskets (silicon
metal, 38\%, and magnet rare earths, 30\%), with shares from 20 to 43\% (Batteries Midstream).
Kazakhstan leads Nuclear Upstream (44\%). In the median basket stage the 37 nodes account for
74\% of world exports (54 to 97\%), so the closed design replaces, by assumption, between 3 and
46\% of observed world supply. Whether the designated nodes could supply that volume, at what
cost and over what horizon, lies outside the model (Section~\ref{subsec:network_limitations}).
Two caveats apply to the benchmarks. Economies that do not report to Comtrade enter the matrix
only through their partners' mirror declarations, and the number of exporters behind each index
differs (34 to 139 world reporters, at most 37 nodes, seven to ten designed producers), so part of
any difference reflects the number of units.

\begin{figure}[htbp]
  \centering
  \includegraphics[width=\textwidth]{network_concentration}
  \caption{Designed versus observed concentration by technology and stage. Panel~(a): Herfindahl
  index of the designed production shares, $\HHI^{\mathrm{D}}_{ts}$ (filled circles coloured by
  stage); of the 2024 basket exports of the 37 network nodes, $\HHI^{\mathrm{N}}_{ts}$ (grey;
  shares within the nodes, the EU-27 as one exporter counting only extra-EU exports; reconciled
  Comtrade matrix, importer report else the exporter's mirror $\times1.06$); and of world exports
  across all reporters, $\HHI^{\mathrm{W}}_{ts}$ (open circles). A grey segment joins the nodes'
  observed value to the design. The vertical line is the calibration target, the median of
  $\HHI^{\mathrm{N}}_{ts}$ over the 18 stages with a basket (0.156). Panel~(b): the largest
  single share in each case. The codes at the right name the world's largest exporter of the
  stage's basket, in bold where it is outside the network; Geothermal and Green Hydrogen have no
  basket and show the design only. Production design ($\lambda=0.78$, caps on). Sources:
  \texttt{allied\_network\_diagnostics.csv}, \texttt{allied\_network\_observed\_hhi.csv} and the
  trade-concentration theme, from UN Comtrade \citep{uncomtrade}.}
  \label{fig:network_concentration}
\end{figure}

\paragraph{Design versus current exporters.} The design does not reproduce the ranking of current
exporters, and is not meant to. In the median basket stage 66\% of designed production goes to
nodes outside the three largest exporting nodes (42 to 92\%); all three leaders are designed
producers in eight of the 18 stages and at least one is in every stage. Across the 37 nodes the
Spearman correlation between designed production share and observed export share has a median of
0.53 over the stages (0.27 to 0.74). Part of the gap is mechanical: export performance enters the
producer score only through the trade category of $\EO$, alongside production, demand, readiness,
costs and the other categories, and through $\ES$. Part is the design's purpose. A program that
maximises the network's score subject to redundancy and a diversification penalty selects
combinations of producers, not the top of a ranking, which is the sense in which the AIB concept
paper asks that an allied design ``optimize combinations, not rankings''. % CITE: neis2026aib
The lowest correlations, in Batteries and Electric Vehicles Upstream (0.27 and 0.29),
illustrate it: in Batteries Upstream the design places Indonesia (0.235) level with Australia, the
largest exporter, and gives Canada and Chile 0.12 each, while South Africa and the Philippines,
the second and third exporters, receive 0.05 and 0.06.

\paragraph{Portfolio caps.} The GDP-log caps \eqref{eq:net_cap} range from two (Gabon and Namibia)
to ten (the United States and the EU) and sum to 218, against the 110 top-5 appearances that the
22 stages create. Eight nodes end at their cap (the United States, the EU, India, Canada,
Indonesia, Thailand, Viet Nam and Malaysia), and eleven have no top-5 appearance. The final pass
bars 49 (node, stage) combinations in 17 stages, most of them for Thailand and the United States
(nine each), Indonesia and Malaysia (eight each). Figure~\ref{fig:network_portfolio} shows what
the caps prevent. Solved without them at the same $\lambda$ (the no-caps run of the calibration
study, whose inputs match the production design to $5\times10^{-16}$), six nodes exceed their cap:
the United States would be a top-5 producer in 17 stages against a cap of ten, Indonesia in 12
(cap seven), India in ten (eight), Thailand and Malaysia in ten (six each) and Canada in eight
(seven). The caps pass these appearances chiefly to Japan and Mexico (from three to six each),
the Philippines (from none to four), South Africa (one to four) and Israel (two to four). They do
not reduce the number of nodes without any role, eight with and without caps, and they cost 1.9\%
of the summed stage objective (at most 5.9\% in a stage; 17 stages lose). They also raise
concentration: without caps the median $\HHI^{\mathrm{D}}_{ts}$ at $\lambda=0.78$ is 0.139,
below the 0.152 of the capped design, so the calibrated $\lambda$ is specific to the capped
configuration. Top-5 membership is decided by a strict rank; four
Upstream stages (Coal, Green Hydrogen, Nuclear and Solar) have a tie at fifth place.

\begin{figure}[htbp]
  \centering
  \includegraphics[width=0.75\textwidth]{network_portfolio}
  \caption{Portfolio caps and realised top-5 appearances. For each of the 37 network nodes (ISO3
  codes; EU = the EU-27 node), ordered by GDP with the largest at the top, the vertical tick is
  the cap $\kappa_i$ of \eqref{eq:net_cap} (GDP-log rule, between two and ten), the filled circle
  the number of stages in which the node is a top-5 producer \eqref{eq:net_topk} in the production
  design, and the open circle the same count in the design solved without portfolio caps and
  otherwise unchanged ($\lambda=0.78$); a segment joins the two where they differ. Counts are over
  the 22 Upstream and Midstream stages. GDP is the latest finite WDI value in current US dollars
  \citep{worldbank2026wdi}, summed over the 27 members for the EU. Source:
  \texttt{allied\_network\_portfolio\_caps.csv}, \texttt{allied\_network\_portfolio\_counts.csv}
  and the no-caps run of the calibration study.}
  \label{fig:network_portfolio}
\end{figure}

\paragraph{Build candidates.} Ranked by the build priority \eqref{eq:net_build} and restricted to
the 22 nodes with an observed development potential, the ten leading candidates are India in
Solar Midstream (priority 0.76), Wind Midstream (0.75), Gas Midstream (0.72) and Nuclear
Midstream (0.69); Indonesia in Electric Vehicles Upstream and Coal Upstream (0.74 each), Batteries
Upstream (0.71), Solar Upstream (0.69) and Wind Upstream (0.68); and Viet Nam in Solar Midstream
(0.67). All ten are already designed producers. This follows from the construction: the priority
multiplies development potential by a weighted mean of the producer score and connectivity, the
quantities on which the program itself selects producers, so the list shows where the design
already relies on an emerging economy rather than where new capacity would add a producer.

\paragraph{Calibration of the diversification weight.} Figure~\ref{fig:network_frontier} and
Table~\ref{tab:network_frontier} document the rule that sets $\lambda$: the smallest value on a
grid at which the median $\HHI^{\mathrm{D}}_{ts}$ over the 22 stages is at most the median export
Herfindahl index of the nodes themselves over the 18 basket stages, 0.1556. The grid has 36
points (0, 0.5, every 0.01 from 0.70 to 1.00, then 1.1, 1.2 and 1.5), each a full re-solve of the
production configuration. The rule picks $\lambda=0.78$, at which the median is 0.1520; the run
at the pick is the stored production design (largest share difference $5\times10^{-16}$). The
pick is knife-edge in three respects. The margin is small: at 0.77 the median is 0.1558 and
misses the target by 0.0002, at 0.76 it is 0.1583. The median is not monotone in $\lambda$: it
rises in eight of the 35 grid steps, among them from 0.1485 at 0.79 to 0.1502 at 0.80, because
interior shares move top-5 membership and the portfolio loop then bars different nodes; no grid
point above the pick falls back above the target. And the pick moved with every revision of the
inputs during the final data corrections, which is why the rule, not the number, is the
specification, and why it is re-applied on the inputs of each full rerun. Every grid run up to
$\lambda=0.91$ solves all 22 stages by the MILP; four runs above it (0.92, 0.93, 0.95 and 0.99)
end Nuclear Upstream on the greedy fallback after the 300-second limit, and the portfolio loop
converges in every run (four to eight passes). Diversification has a price in the score: at the
pick the summed node and edge terms of the 22 programs are 5.8\% below their value at $\lambda=0$,
and 10.6\% below at $\lambda=1.5$, where the median index is 0.095 (10.6 effective producers).
These are differences in an index score, not in resource costs; they rank designs without pricing
them, but they make the trade-off between score and dispersion explicit in the way the AIB concept
paper asks of an allied portfolio objective, as a frontier rather than a single weighted
number. % CITE: neis2026aib

\begin{figure}[htbp]
  \centering
  \includegraphics[width=\textwidth]{network_frontier}
  \caption{Calibration of the diversification weight $\lambda$. Panel~(a): median over the 22
  Upstream and Midstream stages of the Herfindahl index of designed production shares, for each of
  the 36 grid values of $\lambda$ (production configuration: 37 nodes, portfolio caps on, MILP,
  300-second stage limit); the horizontal line is the target, the median export Herfindahl index
  of the nodes over the 18 basket stages (0.1556); the large marker is the pick, $\lambda=0.78$;
  the open circle the same design without portfolio caps; the shaded band the range of
  panel~(b). Panel~(b): the neighbourhood of the pick, $0.70\le\lambda\le0.90$; open markers are
  grid points above the target, filled markers points that meet it, and orange segments the steps
  at which the median rises although $\lambda$ increases. Panel~(c): the summed node and edge terms
  of the 22 stage objectives (without the Herfindahl penalty), in percent of their value at
  $\lambda=0$. Source: calibration study on the inputs of the final rerun of 12 October 2026
  (\texttt{lambda\_study\_g37.csv}, \texttt{lambda\_study\_nc37.csv} and the runs behind them).}
  \label{fig:network_frontier}
\end{figure}

% Generated by scripts/paper/fig_network_frontier.R
\begin{table}[htbp]
\centering
\begin{threeparttable}
\caption{Calibration of the diversification weight: the grid around the pick}
\label{tab:network_frontier}
\footnotesize
\begin{tabular}{@{}l c r r r r r c@{}}
\toprule
$\lambda$ & Caps & Median $\HHI^{\mathrm{D}}$ & Eff.\ producers & Max.\ share & Target $-$ median & MILP stages & Meets \\
\midrule
0 & on & 0.3450 & 2.90 & 0.400 & -0.1894 & 22 & no \\
0.5 & on & 0.2156 & 4.64 & 0.400 & -0.0601 & 22 & no \\
0.74 & on & 0.1620 & 6.17 & 0.335 & -0.0065 & 22 & no \\
0.75 & on & 0.1593 & 6.28 & 0.335 & -0.0037 & 22 & no \\
0.76 & on & 0.1583 & 6.32 & 0.331 & -0.0028 & 22 & no \\
0.77 & on & 0.1558 & 6.42 & 0.325 & -0.0002 & 22 & no \\
\textbf{0.78} & on & 0.1520 & 6.58 & 0.325 & +0.0036 & 22 & yes \\
0.79 & on & 0.1485 & 6.74 & 0.325 & +0.0070 & 22 & yes \\
0.8 & on & 0.1502 & 6.66 & 0.320 & +0.0054 & 22 & yes \\
0.81 & on & 0.1500 & 6.66 & 0.315 & +0.0055 & 22 & yes \\
0.82 & on & 0.1500 & 6.66 & 0.315 & +0.0055 & 22 & yes \\
1 & on & 0.1291 & 7.74 & 0.275 & +0.0264 & 22 & yes \\
1.5 & on & 0.0945 & 10.58 & 0.215 & +0.0611 & 22 & yes \\
\addlinespace
0.78 & off & 0.1393 & 7.18 & 0.312 & +0.0163 & 22 & yes \\
\bottomrule
\end{tabular}
\begin{tablenotes}[flushleft]\footnotesize
\item \emph{Notes:} Production configuration of the network design (37 nodes, the 22 Upstream and Midstream stages, node policy term off, MILP with a 300-second stage limit and the greedy fallback, GDP-log portfolio caps) re-solved at each grid value of $\lambda$. Median $\HHI^{\mathrm{D}}$: median over the 22 stages of the Herfindahl index of designed production shares; Eff.\ producers: median over the stages of its inverse. Target: median over the 18 stages with an HS6 basket of the nodes' observed 2024 export Herfindahl index (the EU-27 node counting extra-EU exports only), 0.1556 (6.43 equivalent exporters). Meets: median $\HHI^{\mathrm{D}}\le$ target; the rule picks the smallest grid value that meets it (bold). The grid runs from 0 to 1.5 with a step of 0.01 between 0.7 and 1; the table shows the points within 0.04 of the pick, the anchors 0, 0.5, 1 and 1.5, and the run without portfolio caps at the pick. Source: lambda study of the final inputs (2026-10-12).
\end{tablenotes}
\end{threeparttable}
\end{table}


\paragraph{Robustness of the design.} Figure~\ref{fig:network_robustness} and
Table~\ref{tab:network_robustness} re-solve the production design with one setting changed at a
time, on inputs assembled exactly as the production script assembles them; the re-run of the
production configuration reproduces the stored design to $5\times10^{-16}$, and the re-runs
without caps and at $\lambda=0.77$ and 0.79 reproduce the corresponding runs of the calibration
study exactly. Two measures compare each variant with the production design stage by stage: the
Jaccard overlap of the producer sets, which treats a producer at the floor like a leading one,
and the share of production reassigned, $\tfrac12\sum_p|x^{\mathrm{alt}}_{p,ts}-x_{p,ts}|$, which
weights each change by its size. No stage has the same producer set in all eight variants, and
none changes in all of them.

The neighbours of the calibrated $\lambda$ leave the design almost unchanged. At $\lambda=0.77$,
21 of the 22 producer sets are identical and 0.5\% of production is reassigned in the mean stage
(Thailand leaves Oil Upstream); at 0.79, 19 sets are identical and 1.1\% is reassigned (Norway
replaces Thailand in Nuclear Upstream, Indonesia enters Batteries Midstream and Israel Electric
Grid Midstream). The pick of $\lambda$ is knife-edge in the sense of the calibration rule, since
the median index crosses the target between these points (0.1558 at 0.77, 0.1520 at 0.78, 0.1485
at 0.79), but the design it selects is not: the step from one grid point to the next moves a
floor share or two, not the division of labour. The demand specification matters as little.
Demand from need alone changes one stage (Singapore enters Coal Midstream at the floor), demand
from size alone changes seven and reassigns 3.4\% of production in the mean stage (at most 15\% in
Geothermal Upstream). Demand enters the choice of producers only through the edge term and the
redundancy constraints, and the S-curved demand weights are rank-compressed (the largest weight
of any node in any stage is 0.054; Section~\ref{subsec:network_demand}), so the producer map is
insensitive to how demand is measured.

The edge weight and the portfolio caps matter more. Without the edge term ($w_E=0$) the mean
Jaccard overlap is 0.84 and 8.4\% of production is reassigned; with an edge term as heavy as the
node term ($w_E=1$) the overlap is 0.83 and 10.8\% is reassigned. The direction is systematic: a
heavier edge term concentrates production on the best-connected producers (median
$\HHI^{\mathrm{D}}_{ts}$ 0.157 at $w_E=1$ against 0.146 at $w_E=0$; largest share 0.355 against
0.295), so the calibrated $\lambda$ is specific to $w_E=0.5$ as well as to the caps. The stages
that move most are Electric Grid Upstream when the edges are dropped (Canada, the EU and the
United States replace the United Kingdom, Japan and Norway) and Electric Vehicles Midstream when
they are doubled (India, Malaysia and Thailand replace the United Kingdom, Viet Nam and South
Africa). Removing the caps is the largest change: five producer sets are unchanged, the mean
overlap is 0.69 (0.29 in Electric Grid Upstream) and 18.5\% of production is reassigned in the
mean stage (46.5\% in Electric Grid Upstream). Without caps the design has 214 producer roles
instead of 201, 49 roles enter and 36 leave; the entrants are chiefly the large and high-scoring
nodes that the caps restrain (Thailand and the United States nine each, Indonesia and Malaysia
eight each) and the leavers the nodes to which the caps pass their appearances (Israel five, the
United Kingdom and Japan four each). The same eight nodes have no role with and without caps.

The MILP also matters. Solved by the greedy heuristic, with the caps on, the design keeps a mean
overlap of 0.84 but only two identical producer sets, is more concentrated (median
$\HHI^{\mathrm{D}}_{ts}$ 0.160, above the calibration target; largest share 0.346) and scores
2.5\% lower in the summed stage objective. On identical eligible sets, without caps, the MILP
objective is at least the greedy one in all 22 stages, by 2.4\% in the median stage and at most
4.0\%, and the greedy design's median index is 0.154 against the MILP's 0.139. The greedy fallback
is therefore a safeguard, not an equivalent: had it been triggered, the design would be more
concentrated and lower-scoring. It is not triggered in the production design or in any variant
solved by the MILP. Across all variants, the least stable stages are Solar Upstream, Electric
Vehicles Midstream, Coal Midstream and Wind Upstream (11 to 14\% of production reassigned in the
mean variant), the most stable Gas Upstream, Geothermal Midstream, Oil Midstream and Batteries
Upstream (2 to 3\%). The variants change one setting at a time; interactions between them, the
node weights of the producer score and the redundancy parameters $m_S$ and $\varepsilon$ are not
explored here. These tests concern the stability of the design to its specification. They do not
test the robustness of the designed network to shocks in the sense of the AIB concept paper's
network-robustness layer (common-mode failure, logistics chokepoints, surge capacity and recovery
time), which the program does not model beyond the redundancy floor of $m_S=2$ suppliers per
consumer. % CITE: neis2026aib

\begin{figure}[htbp]
  \centering
  \includegraphics[width=\textwidth]{network_robustness}
  \caption{Robustness of the network design, by stage and variant. Each column re-solves the 22
  Upstream and Midstream stage programs on the 37 network nodes with one setting of the production
  design changed: no portfolio caps; the greedy heuristic instead of the MILP (caps on); edge
  weight $w_E=0$ or $w_E=1$ instead of 0.5; demand from need only or from size only instead of
  their equal mix; and the neighbouring grid values of the diversification weight, $\lambda=0.77$
  and 0.79, instead of 0.78. Panel~(a): share of the producer set that changes, one minus the
  Jaccard overlap $|A\cap B|/|A\cup B|$ of the stage's producer sets in the variant and the
  production design. Panel~(b): share of the stage's production assigned to different nodes,
  $\tfrac12\sum_p|x^{\mathrm{alt}}_{p,ts}-x_{p,ts}|$. Both are 0 when the stage is unchanged
  (marked by a dot) and 1 when the producer sets are disjoint; darker cells change more.
  Production design: $\lambda=0.78$, GDP-log portfolio caps, $w_E=0.5$, mixed demand, MILP.
  Source: \texttt{scripts/paper/fig\_network\_robustness.R}, which re-runs the design with the
  inputs of the final rerun of 12 October 2026.}
  \label{fig:network_robustness}
\end{figure}

% Generated by scripts/paper/fig_network_robustness.R
\begin{table}[htbp]
\centering
\begin{threeparttable}
\caption{Robustness of the network design to one setting at a time}
\label{tab:network_robustness}
\footnotesize
\begin{tabular}{@{}l r r r r r r r r@{}}
\toprule
 & \multicolumn{4}{c}{Against the production design} & \multicolumn{4}{c}{Design of the variant} \\
\cmidrule(lr){2-5}\cmidrule(l){6-9}
Variant & Jaccard & Min. & \makecell[r]{Identical\\stages} & \makecell[r]{Re-\\assigned} & \makecell[r]{Median\\$\HHI^{\mathrm{D}}$} & \makecell[r]{Pro-\\ducers} & \makecell[r]{Max.\\share} & \makecell[r]{MILP\\stages} \\
\midrule
Production design & -- & -- & -- & -- & 0.152 & 9 & 0.325 & 22 \\
\addlinespace
No portfolio caps & 0.692 & 0.294 & 5 & 0.185 & 0.139 &   10 & 0.312 & 22 \\
Greedy heuristic & 0.836 & 0.600 & 2 & 0.086 & 0.160 &    9 & 0.346 & 0 \\
Edge weight $w_E=0$ & 0.840 & 0.538 & 6 & 0.084 & 0.146 &    9 & 0.295 & 22 \\
Edge weight $w_E=1$ & 0.825 & 0.455 & 7 & 0.108 & 0.157 &    9 & 0.355 & 22 \\
Demand from need only & 0.995 & 0.900 & 21 & 0.010 & 0.154 &    9 & 0.325 & 22 \\
Demand from size only & 0.941 & 0.600 & 15 & 0.034 & 0.159 &    9 & 0.325 & 22 \\
$\lambda=0.77$ & 0.995 & 0.889 & 21 & 0.005 & 0.156 &    9 & 0.325 & 22 \\
$\lambda=0.79$ & 0.980 & 0.750 & 19 & 0.011 & 0.149 &    9 & 0.325 & 22 \\
\bottomrule
\end{tabular}
\begin{tablenotes}[flushleft]\footnotesize
\item \emph{Notes:} Production design: $\lambda=0.78$, GDP-log portfolio caps, edge weight $w_E=0.5$, demand from the equal mix of need and size, MILP. Each other row re-solves the 22 Upstream and Midstream stage programs on the 37 network nodes with one setting of the production design changed. Jaccard: mean across stages of the Jaccard overlap $|A\cap B|/|A\cup B|$ of the stage's producer set with that of the production design; Min.: its smallest value; Identical stages: stages with the same producer set. Reassigned: mean across stages of the share of production assigned to different nodes, $\tfrac12\sum_p|x^{\mathrm{alt}}_{p,ts}-x_{p,ts}|$ (0 when the shares are unchanged, 1 when the producer sets are disjoint). Median $\HHI^{\mathrm{D}}$: median across stages of the Herfindahl index of designed shares. Producers: median number of producers per stage. Max. share: largest designed share in any stage. MILP stages: stages solved by the MILP (the greedy row solves all stages heuristically). Source: \texttt{scripts/paper/fig\_network\_robustness.R}.
\end{tablenotes}
\end{threeparttable}
\end{table}


